\chapter{مفاهيم أساسية}

%\section{العمليةﻹ الثنائية}

\begin{definition}{العملية الثنائية}{\cite{ref1}}
	\addcontentsline{toc}{section}{العملية الثنائية}
	لتكن $S$ مجموعة غير خالية فأن اي تطبيق من 
	$S \times S$ الى $S$ يسمى عملية ثنائية ويرمز له بإحدى الرموز التالية
	($*, \#, \$, \circ, + , \div, \cap, \cup, \Delta$) 
\end{definition}

\begin{note}
	تكون العملية $*$ مغلقة على المجموعة $S$ اذا تحقق الشرط التالي
	\[
	\forall a, b\in S \implies a * b \in S
	\]
\end{note}

\begin{example}
	عملية الجمع $+$ تكون ثنائية على مجموعة الاعداد الطبيعية $\N$ 
	\[
	\forall a, b \in \N \implies a + b \in \N
	\]
\end{example}

%\section{النظام الرياضي}

\begin{definition}{النظام الرياضي}{\cite{ref1}}
	النظامُ الرياضي هو زوجٌ مرتب $(S, *)$ حيث إنَّ $S$ مجموعة غير خالية و$*$ عملية ثنائية معرفة على $S$.
\end{definition}

\begin{definition}{شبه الزمرة}{\cite{ref1}}
	شبهُ الزمرة هو نظامٌ رياضي $(S, *)$ حيث إنَّ $S$ مجموعة غير خالية و$*$ عملية ثنائية معرفة على $S$ وتحقق خاصية التجميع، أي:
	\[
	a * (b * c) = (a * b) * c \quad \forall a, b, c \in S.
	\]
\end{definition}

\begin{example}
	النظام الرياضي $(\N, +)$ يكون شبه زمرة لأن
	\begin{enumerate}[nosep]
		\item عملية الجمع $+$ مغلقة على $\N$.
		\item عملية الجمع تكون تجميعية:
		\[
		\forall a, b, c\in \N : a + (b + c) = (a + b) + c 
		\]
	\end{enumerate}
\end{example}

%\section{الزمرة}

\begin{definition}{الزمرة}{\cite{ref2}}
	\addcontentsline{toc}{section}{الزمرة}
	يكون النظام الرياضي $(G, *)$ حيث $G$ مجموعة غير خالية و $*$ عملية ثنائية على $G$ زمرة (group) إذا تحقق ما يلي:
	\begin{enumerate}
		\item المجموعة $G$ مغلقة تحت العملية $*$، أي أن:
		\[
		\forall a, b \in G \Rightarrow a * b \in G
		\]
		
		\item العملية $*$ تحقق خاصية التجميع، أي أن:
		\[
		\forall a, b, c \in G \Rightarrow a * (b * c) = (a * b) * c
		\]
		
		\item يوجد عنصر محايد $e \in G$ بحيث:
		\[
		a * e = e * a = a \quad \forall a \in G
		\]
		
		\item لكل عنصر $a \in G$ يوجد عنصر معاكس $b \in G$ بحيث:
		\[
		a * b = b * a = e
		\]	
	\end{enumerate}
	
\end{definition}

\begin{note}
\begin{tasks*}
	\item 	يسمى العنصر $e$ في التعريف اعلاه بالعنصر المحايد 
	(identity element).
	
	\item يسمى العنصر $b$ في التعريف اعلاه بالعنصر النظير للعنصر $a$ (\en{inverse of $a$}) ويرمز عادة بــــــ $a^{-1}$.
	
	\item في العادة اذا كانت $(G, *)$ زمرة فإننا نستخدم عملية الضرب بدلاً من $*$ ونكتب $ab$ بدلاً من $a*b$. كما نكتب أحياناً $G$ بدلاً من 
	$(G, *)$.
\end{tasks*}
\end{note}

\begin{example}
	لتكن المجموعة \( G = \mathbb{Z} \) مع العملية \( + \). نتحقق من أنها زمرة كما يلي:
	\begin{enumerate}[label=(\alph*)]
		\item \textbf{الإغلاق (Closure):}  
		إذا كان \( a, b \in \mathbb{Z} \) فإن \( a + b \in \mathbb{Z} \)، إذن العملية مغلقة.
		\item \textbf{التجميعية (Associativity):}  
		لكل \( a, b, c \in \mathbb{Z} \) لدينا:
		\[
		(a + b) + c = a + (b + c)
		\]
		إذن العملية تجميعية.
		\item \textbf{العنصر المحايد (Identity):}  
		يوجد عنصر \( 0 \in \mathbb{Z} \) بحيث:
		\[
		a + 0 = a \quad \forall a \in \mathbb{Z}
		\]
		\item \textbf{المعكوس (Inverse):}  
		لكل \( a \in \mathbb{Z} \) يوجد عنصر \( -a \in \mathbb{Z} \) بحيث:
		\[
		a + (-a) = 0
		\]
	\end{enumerate}
	إذن، تحقق جميع الشروط، وبالتالي فإن \( (\mathbb{Z}, +) \) زمرة.
\end{example}

\begin{theorem}[\cite{ref2}]
	اذا كانت $(G, *)$ زمرة فإن:
	\begin{enumerate}[label=(\alph*)]
		\item العنصر المحايد $e$ يكون وحيد.
		\item لكل $a\in G$ العنصر النظير $b$ يكون وحيد.
	\end{enumerate}
\end{theorem}
\begin{proof}
	\begin{tasks*}
		\item لنفرض أن $e_1 \in G$ عنصر محايد آخر يحقق $a * e_1 = e_1 * a = a, \forall a\in G$. عندئذٍ :\linebreak $e_1 = e_1 * e = e*e_1$.
		\item لنفرض $c\in G$ عنصر نظير آخر للعنصر $a\in G$ اي ان : 
		$\forall a \in G  : a * c= c*a = e$. عندئذٍ:\linebreak
		$b = b*e = b*(a*c) = (b*a) * c = e*c = c$.
	\end{tasks*}
\end{proof}

\begin{note}
	تسمى الزمرة $G$ زمرة ابدالية اذا حققت الشرط التالي
	\[
	ab = ba \forall a, b\in G
	\]
\end{note}

\begin{example}
	كل من
	$\Z, \Q, \R, \C$ زمرة ابدالية مع عملية الجمع حيث العنصر $0$ هو العنصر المحايد ونظير\linebreak العنصر $a$ هو $-a$.
\end{example}

%\section{الزمرة الجزئية}
\begin{definition}{الزمرة الجزئية}{\cite{ref3}}\\
	\addcontentsline{toc}{section}{الزمرة الجزئية}
	لتكن $(G, *)$ زمرة، ولتكن $H$ مجموعة جزئية غير خالية من $G$. تُسمى $H$ زمرةً جزئية من $G$ إذا تحققت الشروط التالية:
	\begin{enumerate}
		\item المجموعة $H$ مغلقة تحت العملية $*$، أي أن:
		\[
		\forall a, b \in H \Rightarrow a * b \in H
		\]
		
		\item لكل عنصر $a \in H$ يوجد معكوسه في $H$، أي أن:
		\[
		\forall a \in H \Rightarrow a^{-1} \in H
		\]
	\end{enumerate}
\end{definition}

\begin{example}
	لتكن $(\mathbb{Z}, +)$ زمرة الأعداد الصحيحة تحت عملية الجمع، ولتكن
	\[
	H = 2\mathbb{Z} = \{ \dots, -4, -2, 0, 2, 4, \dots \}
	\]
	نبيّن أن $H$ زمرة جزئية من $\mathbb{Z}$:
	
	\begin{enumerate}
		\item الانغلاق: إذا كان $a, b \in H$ فإن $a = 2m$ و$b = 2n$ لبعض $m, n \in \mathbb{Z}$، وبالتالي:
		\[
		a + b = 2m + 2n = 2(m+n) \in H
		\]
		
		\item المعكوس: إذا كان $a \in H$ فإن $a = 2m$ لبعض $m \in \mathbb{Z}$، وبالتالي:
		\[
		-a = -2m = 2(-m) \in H
		\]
	\end{enumerate}
	إذن $H$ زمرة جزئية من $(\mathbb{Z}, +)$.
\end{example}

\begin{definition}{الزمرة الجئية الناظمية}{\cite{ref3}}
	لتكن $(G, *)$ زمرة و $H$ زمرة جزئية من $G$. تُسمى $H$ زمرةً ناظمية في $G$ إذا تحقق الشرط:
	\[
	a * H = H * a \quad \forall a \in G
	\]
	حيث:
	\[
	a * H = \{a * h \mid h \in H\}, \quad H * a = \{h * a \mid h \in H\}.
	\]
	ويرمز لها بـ $H \trianglelefteq G$.
\end{definition}

\begin{example}
	كل زمرة جزئية من زمرة تبديلية تكون زمرة ناظمية. مثلاً:
	\[
	(\mathbb{Z}, +)
	\]
	فكل زمرة جزئية مثل $n\mathbb{Z}$ هي زمرة ناظمية في $\mathbb{Z}$.
\end{example}

\begin{example}
	في زمرة التباديل $S_3$، تكون المجموعة:
	\[
	A_3 = \{ e, (123), (132) \}
	\]
	زمرة ناظمية في $S_3$.
\end{example}

\begin{example}
	في الزمرة $S_n$، فإن الزمرة التناوبية $A_n$ هي زمرة ناظمية في $S_n$ لكل $n \ge 2$.
\end{example}

%\section{التباديل \en{\textit{{S\textsubscript{n}}}}}

\begin{definition}{التباديل}{\cite{ref4}}
	\addcontentsline{toc}{section}{التباديل}
	لتكن $A = \{1, 2, \dots, n\}$ مجموعة منتهية. يُسمى كل تقابل (تطبيق تقابلي) من $A$ إلى نفسها تبديلاً. وتُسمى مجموعة جميع التباديل على $A$ بالزمرة المتماثلة من الدرجة $n$، ويرمز لها بـ $S_n$.
\end{definition}

تكون العملية المعرفة على $S_n$ هي تركيب التباديل، حيث إن تركيب تبديلين يعطي تبديلاً أيضاً، وبالتالي فإن $(S_n, \circ)$ زمرة.

\begin{example}
	في الحالة $n = 2$، لدينا:
	\[
	S_2 = \{ e, (12) \}
	\]
	حيث $e$ هو التبديل المحايد و$(12)$ يبدل بين $1$ و$2$.
\end{example}

\begin{example}
	في الحالة $n = 3$، لدينا:
	\[
	S_3 = \{ e, (12), (13), (23), (123), (132) \}
	\]
	حيث:
	\begin{itemize}
		\item التباديل $(12), (13), (23)$ تُسمى تبديلات ثنائية.
		\item التباديل $(123), (132)$ تُسمى دورات ثلاثية.
	\end{itemize}
\end{example}

\begin{example}
	ليكن $\sigma \in S_3$ معرفاً بـ:
	\[
	\sigma =
	\begin{pmatrix}
		1 & 2 & 3 \\
		2 & 3 & 1
	\end{pmatrix}
	\]
	فإن هذا التبديل يمكن كتابته على صورة دورة:
	\[
	\sigma = (123)
	\]
\end{example}

\begin{note}
	زمرة التباديل $S_n$ تكون غير إبدالية عندما يكون $n \ge 3$، أي أنه يوجد تبديلان $\sigma, \tau \in S_n$ بحيث:
	\[
	\sigma \circ \tau \ne \tau \circ \sigma.
	\]
	أما في الحالة $n = 1$ أو $n = 2$ فإن الزمرة تكون إبدالية.
\end{note}

%\section{علاقات التكافؤ}

\begin{definition}{العلاقة}{\cite{ref1}}
	لتكن $A$ مجموعة غير خالية. تُعرَّف العلاقة على $A$ بأنها أي مجموعة جزئية من $A \times A$.
\end{definition}

\begin{definition}{علاقة التكافؤ}{\cite{ref1}}
	تسمى العلاقة $R$ على مجموعة $A$ علاقة تكافؤ إذا حققت الشروط التالية:
	\begin{enumerate}
		\item الانعكاسية: $\forall a \in A \Rightarrow aRa$
		\item التناظر: $\forall a,b \in A \Rightarrow aRb \Rightarrow bRa$
		\item التعدي: $\forall a,b,c \in A \Rightarrow aRb \wedge bRc \Rightarrow aRc$
	\end{enumerate}
\end{definition}

\begin{definition}{صفوف التكافؤ}{\cite{ref1}}
	لتكن $R$ علاقة تكافؤ على مجموعة $A$. يُسمى صف التكافؤ للعنصر $a \in A$ هو:
	\[
	[a] = \{ x \in A \mid xRa \}
	\]
	وتسمى مجموعة جميع صفوف التكافؤ بـ مجموعة القسمة، وترمز:
	\[
	A / R = \{ [a] \mid a \in A \}.
	\]
\end{definition}

\begin{example}
	لتكن $A = \mathbb{Z}$ ونعرّف العلاقة $R$ كما يلي:
	\[
	aRb \iff 3 \mid (a - b)
	\]
	أي أن $a-b$ يقبل القسمة على $3$.
	
	نثبت أن $R$ علاقة تكافؤ:
	
	\begin{enumerate}
		\item الانعكاسية: لكل $a \in \mathbb{Z}$ لدينا
		\[
		a-a = 0 \Rightarrow 3 \mid 0 \Rightarrow aRa
		\]
		
		\item التناظر: إذا كان $aRb$ فإن
		\[
		3 \mid (a-b) \Rightarrow 3 \mid -(a-b) = (b-a) \Rightarrow bRa
		\]
		
		\item التعدي: إذا كان $aRb$ و $bRc$ فإن
		\[
		3 \mid (a-b), \quad 3 \mid (b-c)
		\]
		وبالتالي:
		\[
		(a-c) = (a-b) + (b-c) \Rightarrow 3 \mid (a-c) \Rightarrow aRc
		\]
	\end{enumerate}
	إذن $R$ علاقة تكافؤ على $\mathbb{Z}$.
	
	نحسب صفوف التكافؤ:
	
	\begin{enumerate}
		\item صف العدد $0$:
		\[
		[0] = \{ \dots, -6, -3, 0, 3, 6, \dots \}
		\]
		
		\item صف العدد $1$:
		\[
		[1] = \{ \dots, -5, -2, 1, 4, 7, \dots \}
		\]
		
		\item صف العدد $2$:
		\[
		[2] = \{ \dots, -4, -1, 2, 5, 8, \dots \}
		\]
	\end{enumerate}
	وبالتالي:
	\[
	\mathbb{Z}/R = \{ [0], [1], [2] \}.
	\]
\end{example}

%\section{التشاكل الزمري}

\begin{definition}{التشاكل الزمري}{\cite{ref3}}
	\addcontentsline{toc}{section}{التشاكل الزمري}
	لتكن $(G, *)$ و $(H, \circ)$ زمرتين. نقول إن الدالة $\varphi : G \to H$ هي تشاكل زُمَري إذا حققت الشرط
	 \linebreak$\varphi(a * b) = \varphi(a) \circ \varphi(b)$ لكل $a, b \in G$.
\end{definition}


\begin{example}
	لتكن الدالة $f : (\mathbb{Z}, +) \to (\mathbb{R}^*, \cdot)$ معرفة بـ $f(x) = e^x$. عندئذٍ تكون $f$ تماثلًا زمريًا لأن
	\linebreak $f(x+y) = e^{x+y} = e^x \cdot e^y = f(x)f(y)$ لكل $x,y \in \mathbb{Z}$.
\end{example}

\begin{definition}{زمرة الاعداد الصحيحة معيار \en{\textit{n}}}{\cite{ref3}}
	لتكن $n \in \mathbb{N}$ حيث $n \ge 1$. نعرّف $\mathbb{Z}_n = \{ [0], [1], \dots, [n-1] \}$ حيث \linebreak$[a] = \{ b \in \mathbb{Z} : a \equiv b \pmod{n} \}$، وتُعرَّف العمليات بـ $[a] + [b] = [a+b]$ و$[a][b] = [ab]$.
\end{definition}
\noindent\textbf{مثال}
\[
\Z_6 = \{[0], [1], [2], [3], [4], [5]\}
\]

\begin{definition}{نواة التشاكل}{\cite{ref6}}
	لتكن $\varphi : G \to H$ تشاكلًا زُمريًا بين الزمرتين $(G, *)$ و $(H, \circ)$.
	نُعرّف نواة التشاكل بأنها المجموعة
	\[
	\ker(\varphi) = \{\, a \in G \mid \varphi(a) = e_H \,\}
	\]
	حيث $e_H$ هو العنصر المحايد في الزمرة $H$.
\end{definition}


\begin{example}
	لتكن الدالة $f : (\mathbb{R}^+, \cdot) \to (\mathbb{R}, +)$ معرفة بـ
	\[
	f(x) = \ln(x).
	\]
	عندئذٍ تكون $f$ تماثلًا زُمريًا لأن
	\[
	f(xy) = \ln(xy) = \ln(x) + \ln(y) = f(x) + f(y)
	\]
	لكل $x,y \in \mathbb{R}^+$.
	وبالتالي فإن نواة هذا التشاكل هي:
	\[
	\ker(f) = \{x \in \mathbb{R}^+ \mid \ln(x) = 0\} = \{1\}.
	\]
\end{example}



\begin{definition}{المجال المقابل}{}
	لتكن $f: A \to B$ دالة. المجموعة $B$ تُسمّى المجال المقابل (Codomain) للدالة $f$.
\end{definition}

\begin{definition}{المدى}{}
	لتكن $f: A \to B$ دالة. يُعرّف مدى الدالة $f$ بأنه مجموعة صور عناصر $A$، أي:
	\[
	\operatorname{Im}(f) = \{ f(x) \mid x \in A \}.
	\]
\end{definition}

\begin{definition}{الدالة المتباينة}{}
	لتكن $f: A \to B$ دالة. نقول إن $f$ دالة متباينة إذا كان
	\[
	f(x_1) = f(x_2) \Rightarrow x_1 = x_2
	\]
	لكل $x_1, x_2 \in A$.
\end{definition}

\begin{definition}{الدالة الشاملة}{}
	لتكن $f: A \to B$ دالة. نقول إن $f$ دالة شاملة إذا كان لكل عنصر $y \in B$ يوجد عنصر $x \in A$ بحيث
	\[
	f(x) = y.
	\]
\end{definition}

\begin{definition}{الدالة التقابل}{}
	لتكن $f: A \to B$ دالة. نقول إن $f$ دالة تقابل إذا كانت متباينة وشاملة في آنٍ واحد.
\end{definition}

\begin{definition}{زمرة التماثل}{}
	\addcontentsline{toc}{section}{زمرة التماثل}
عناصرها تكون عبارة عن دوال وخواص هذه الدوال (دالة متقابلة) من المجموعة $A$ الى نفسها ونرمز لها بالرمز 
$\operatorname{symm}(A)$

العمليات التي سوف نستخدمها في زمرة التماثل هي عملية تركيب الدوال
\end{definition} 

\begin{example}
	لتكن $A = \{a, b, c\}$. نريد إيجاد زمرة التماثل $\operatorname{symm}(A)$.
	
	زمرة التماثل هي مجموعة جميع التبديلات (الدوال التقابلية من $A$ إلى نفسها)، وهي:
	\[
	\operatorname{symm}(A) =
	\{ 
	f_1, f_2, f_3, f_4, f_5, f_6
	\}
	\]
	حيث:
	\begin{align*}
		f_1 =
		\begin{pmatrix}
			a & b & c \\
			a & b & c
		\end{pmatrix} ,
		f_2 =
		\begin{pmatrix}
			a & b & c \\
			b & a & c
		\end{pmatrix} ,
		f_3 =
		\begin{pmatrix}
			a & b & c \\
			c & b & a
		\end{pmatrix} 
	\end{align*}
	\[
	f_4 =
	\begin{pmatrix}
		a & b & c \\
		a & c & b
	\end{pmatrix} ,
	f_5 =
	\begin{pmatrix}
		a & b & c \\
		b & c & a
	\end{pmatrix} ,
	f_6 =
	\begin{pmatrix}
		a & b & c \\
		c & a & b
	\end{pmatrix}
	\]
	الآن نجري عملية التركيب بين عنصرين، لنأخذ:
	\[
	f_2 =
	\begin{pmatrix}
		a & b & c \\
		b & a & c
	\end{pmatrix},
	\quad
	f_4 =
	\begin{pmatrix}
		a & b & c \\
		a & c & b
	\end{pmatrix}
	\]
	نحسب $f_2 \circ f_4$ (نطبق $f_4$ أولاً ثم $f_2$):
	\[
	f_2 \circ f_4 =
	\begin{pmatrix}
		a & b & c \\
		f_2(f_4(a)) & f_2(f_4(b)) & f_2(f_4(c))
	\end{pmatrix}
	\]
	نحسب الصور:
	\[
	f_4(a)=a \Rightarrow f_2(a)=b,\quad
	f_4(b)=c \Rightarrow f_2(c)=c,\quad
	f_4(c)=b \Rightarrow f_2(b)=a
	\]
	
	إذن:
	\[
	f_2 \circ f_4 =
	\begin{pmatrix}
		a & b & c \\
		b & c & a
	\end{pmatrix}
	= f_5
	\]
	لنجد معكوس $f_5$:
	\[
	f_5 =
	\begin{pmatrix}
		a & b & c \\
		b & c & a
	\end{pmatrix}
	\]
	نعكس الأسهم:
	\[
	a \mapsto b,\; b \mapsto c,\; c \mapsto a
	\Rightarrow
	b \mapsto a,\; c \mapsto b,\; a \mapsto c
	\]
	إذن:
	\[
	f_5^{-1} =
	\begin{pmatrix}
		a & b & c \\
		c & a & b
	\end{pmatrix}
	= f_6
	\]
\end{example}

\begin{theorem}[(مبرهنة كيلي) \cite{ref3}]
	اذا كانت $G$ زمرة فأن $G$ تماثل زمرة جزئية من $S_G$.
\end{theorem}
